\documentclass[../main]{subfiles}
\begin{document}

\section{Lubricantes}

\img{images/08/logo}{Lubricantes}


\info{Definición de Lubricación}{Es el proceso o técnica empleada para reducir el rozamiento entre dos superficies muy próximas y en movimiento una respecto de la otra, interponiendo para ello una sustancia entre ambas denominada lubricante.}

La lubricación también puede describir fenómenos donde tal reducción del rozamiento ocurra sin intervención humana.

Una adecuada lubricación permite un funcionamiento continuo y suave de los equipos mecánicos, con un ligero desgaste, y sin excesivo estrés o ataque a las partes móviles (cojinetes y engranajes). Cuando falla la lubricación, los metales y otros materiales pueden rozar, desgastarse, perdiendo eficacia, causan desprendimiento de calor, pudiendo llegar a destruirse unos a otros, causando daños irreparables, y fallo general.

\img{images/08/lubricando}{Lubricante en motores.}

\info{Definición de Lubricante}{Un lubricante es una sustancia que, colocada entre dos piezas móviles forma una capa que impide su contacto, permitiendo su movimiento incluso a elevadas temperaturas y presiones. El lubricante es una sustancia (gaseosa, líquida o sólida) que reemplaza una fricción entre dos piezas en movimiento relativo por la fricción interna de sus moléculas, que es mucho menor.}

\subsection{Descripción}


El lubricante es una sustancia que introducida entre dos superficies móviles reduce la fricción entre ellas, facilitando el movimiento y reduciendo el desgaste. Además el lubricante cumple varias otras funciones dentro de una máquina o motor, entre ellas disuelve y transporta al filtro las partículas fruto de la combustión y el desgaste, distribuye la temperatura desde la parte inferior a la superior actuando como un refrigerante, evita la corrosión por óxido en las partes del motor o máquina, evita la condensación de vapor de agua y sella actuando como una junta determinados componentes. 



La propiedad del lubricante de reducir la fricción entre partes se conoce como "Lubricación", y \emph{la ciencia que la estudia es la tribología.}

Un lubricante se compone de una base, que puede ser mineral o sintética y un conjunto de aditivos que le confieren sus propiedades y determinan sus características. Cuanto mejor sea la base menos aditivos necesitará, sin embargo se necesita una perfecta comunión entre estos aditivos y la base, pues sin ellos la base tendría condiciones de lubricación mínimas. Los lubricantes se clasifican según su base como, minerales, vegetales o sintéticos.

\subsection{Tipos de lubricantes por composición y presentación }
\begin{itemize}
     
\item Líquidos: de origen mineral o vegetal. Son necesarios para la lubricación hidrodinámica y son usados comúnmente en la industria, motores y como lubricantes de perforación.
\item Semisólidos: son las denominadas "grasas". Su composición puede ser mineral o vegetal y frecuentemente son combinadas con muchos tipos de lubricantes sólidos como el Grafito, Molibdeno o Litio.
\item Sólidos: es un tipo de material que ofrece mínima resistencia molecular interna por lo que por su composición ofrece óptimas condiciones de lubricación sin necesidad de un aporte lubricante líquido o semisólido. El más común es el grafito aunque la industria está avanzando en investigación en materiales de origen metálico.
\end{itemize}

\subsection{Tipos de lubricantes por su naturaleza }
\img{images/08/tipos}{Lubricante según la naturaleza de su base.}
\begin{itemize}
\item Minerales: son los aceites provenientes del refinado del petróleo.
\item Sintéticos: son creados de forma sintética y no tienen origen natural. Tienen mayor resistencia térmica y mejores propiedades anti-desgaste.
\item  Semi-sintéticos: son mezclas de ambos aceites lo cual le da propiedades diferentes a las que poseen cada uno en forma individual.
\end{itemize}



\subsection{Lubricante mineral}

Es el más usado y barato de las bases parafínicas o nafténicas. Se obtiene tras la destilación del petróleo crudo después del gasóleo y antes que el alquitrán, comprendiendo un 50\% del total del crudo, este hecho así como su precio hacen que sea el más utilizado.

Existen dos tipos de lubricantes minerales clasificados por la industria, grupo 1 y grupo 2 atendiendo a razones de calidad y pureza predominando el grupo 1. Es una base de bajo índice de viscosidad natural (SAE 15) por lo que necesita de gran cantidad de aditivos para ofrecer unas buenas condiciones de lubricación. El origen del lubricante mineral por lo tanto es orgánico, puesto que proviene del petróleo.

Los lubricantes minerales obtenidos por destilación del petróleo son fuertemente aditivados para poder:
\begin{enumerate}
    
\item Soportar diversas condiciones de trabajo.
\item Lubricar a altas temperaturas.
\item Permanecer estable en un amplio rango de temperatura.
\item Tener la capacidad de mezclarse adecuadamente con el refrigerante (visibilidad).
\item Tener un índice de viscosidad alto.
\item Tener higroscopicidad definida, (capacidad de retener humedad).
\end{enumerate}

\subsection{Lubricante sintético}
Es una base artificial y por lo tanto del orden de 3 a 5 veces más costosa de producir que la base mineral. Se crea en laboratorio y puede o no provenir del petróleo. Poseen unas excelentes propiedades de estabilidad térmica y resistencia a la oxidación, así como un elevado índice de viscosidad natural (SAE 30). Poseen un coeficiente de tracción muy bajo, con lo cual se obtiene una buena reducción en el consumo de energía.

    
\subsection{Aditivos de los lubricantes}

La base de un lubricante por sí sola no ofrece toda la protección que necesita un motor o componente industrial, por lo que en la fabricación del lubricante se añaden ciertos aditivos atendiendo a las necesidades del fabricante del motor (Homologación o Nivel autorizado) o al uso al que va a ser destinado el lubricante en cuestión.

Los aditivos usados en el lubricante son:
\begin{enumerate}
    
\item Antioxidantes: Retrasan el envejecimiento prematuro del lubricante.

\item Antidesgaste Extrema Presión (EP): Forman una fina película en las paredes a lubricar. Se emplean mucho en lubricación por barboteo (Cajas de cambio y diferenciales)

\item Antiespumantes: Evitan la oxigenación del lubricante por cavitación reduciendo la tensión superficial y así impiden la formación de burbujas que llevarían aire al circuito de lubricación.

\item Antiherrumbre: Evita la formación de óxido en las paredes metálicas internas del motor y la condensación de vapor de agua.

\item Detergentes: Son los encargados de arrancar los depósitos de suciedad fruto de la combustión.

\item Dispersantes: Son los encargados de transportar la suciedad arrancada por los aditivos detergentes hasta el filtro o cárter del motor.

\item Espesantes: Es un compuesto de polímeros que por acción de la temperatura aumentan de tamaño aumentando la viscosidad del lubricante para que siga proporcionando una presión constante de lubricación.

\item Diluyentes: Es un aditivo que reduce los microcristales de cera para que fluya el lubricante a bajas temperaturas.
 
\end{enumerate}

\subsection{Clasificaciones de aceites lubricantes}

Existen diversos tipos de clasificaciones de aceites lubricantes según el ámbito geográfico, sus propiedades y el fabricante de la máquina a lubricar.

Según el ámbito geográfico podemos encontrar la clasificación americana API (American Petroleum Institute), la clasificación japonesa JASO (Japanese Automotive Standard Organization) y la europea ACEA (Asociación de Constructores Europeos Asociados).

Sus propiedades se clasifican según la norma SAE (Society of Automotive Engineers) que básicamente separa el comportamiento del lubricante a temperatura de -18 °C y la define con una letra W proveniente del inglés "Winter" (Invierno) y otra letra que define el comportamiento del lubricante a temperatura de trabajo 95 °C-105 °C. La tabla SAE hace referencia a las tolerancias que debe satisfacer el lubricante tanto a temperatura ambiente como a temperatura de trabajo, siempre teniendo en cuenta la temperatura interna del motor y adicionalmente la temperatura exterior que si bien influye algo en el comportamiento no es la más importante a la hora de elegir un lubricante adecuado.

\img{images/08/viscocidades}{Lubricantes Multigrados.}

Según el fabricante del motor o componente a lubricar existen normativas de fabricante con diversas nomenclaturas tipo VW505.01, GM Dexos2, Dexron III, MB229.51, LL-01, entre otras. Los fabricantes de motores y componentes conocen al detalle su producto y son conscientes de la importancia de un lubricante adecuado y de las consecuencias de no utilizar un lubricante adecuado. Para "protegerse" y distinguirse de sus competidores hace ya muchos años los fabricantes comenzaron a definir estándares de fabricación de los lubricantes aptos para sus productos. Son las llamadas "Homologaciones del fabricante", que es la prueba de que el lubricante ha sido testado por el fabricante en el motor y por ello expide su correspondiente certificado de homologación.

Los acuerdos comerciales de los responsables de cada marca de vehículos, motores o componentes en cada país con las diferentes empresas petroleras hacen que estas últimas presenten los certificados de homologación exclusivamente de los fabricantes con los que ha llegado a acuerdo dificultando la diagnosis del lubricante adecuado para cada vehículo.

En todo caso cabe destacar que usando un lubricante con la homologación del fabricante de la máquina o vehículo las demás clasificaciones son complementarias. Hay más de 72 homologaciones en el sector de lubricación automotriz debido a la reciente incorporación de filtros de partículas y sistemas anticontaminación y hay fabricantes que disponen de varias normativas de homologación.
\subsection{Grasa (lubricante)}

\img{images/08/grasa}{Grasa roja para cojinetes de ruedas para uso en vehículos.}

La grasa es un lubricante semisólido. Por lo general la grasa consiste de un jabón emulsionado con aceite mineral o vegetal. La característica distintiva de las grasa es que tienen una muy elevada viscosidad inicial, la cual al aplicar un esfuerzo de corte, disminuye dando el efecto de un rodamiento lubricado por aceite con aproximadamente la misma viscosidad del aceite base usado en la formulación de la grasa. Este cambio en la viscosidad es denominado adelgazamiento por cizalladura. A veces la grasa es utilizada para describir materiales lubricantes que son sólidos blandos o líquidos con muy elevada viscosidad, pero estos materiales no presentan las propiedades de adelgazamiento por cizalladura de la grasa convencional. Por ejemplo, la vaselina por lo general no es clasificada como grasa.

Las grasas se aplican a mecanismos que pueden ser lubricados con poca frecuencia y donde la lubricación mediante un aceite no permaneceria en el sitio requerido. Además también actúan como selladores previniendo el ingreso de agua y materiales incompresibles. Los cojinetes lubricados con grasa poseen mejores características de fricción a causa de su elevada viscosidad.



\newpage
\actividad{Leer y responder en base al texto.}
\begin{enumerate}
    \item ¿Qué estudia la tribología? 
    \item ¿Qué entiende por lubricante?
    \item ¿Que tipos de lubricante hay según su presentación, composición y naturaleza?
    \item ¿Que características tiene un lubricante mineral?
    \item ¿Cuales Son las características que tiene un lubricante sintético diferentes al mineral?
\end{enumerate}
\actividad{Investigar}
\begin{enumerate}
    \item ¿Qué es la viscosidad y como se mide?
    \item Vamos a comparar precios, de un lubricante sintético 20W50 (normalmente para motos) con uno de iguales características pero mineral. ¿Qué diferencia hay?
    
    \web{https://listado.mercadolibre.com.ar/lubricante-mineral-20w50}{Lubricantes minerales}
    \web{https://listado.mercadolibre.com.ar/lubricante-sintetico-20w50}{Lubricantes Sintéticos}
    
    
\end{enumerate}
\end{document}
